% MOSIP Decode 2026 - PS1 Project Handbook
%
% One typeface throughout: Latin Modern Mono for headings, body, tables and commands
% alike. To switch the whole document to another single family, change the three
% font lines marked FONT below; nothing else sets a font.
%
% Build:  cd docs/handbook && pdflatex handbook.tex && pdflatex handbook.tex
% Uses only packages present in a stock BasicTeX install (no enumitem, no titlesec).

\documentclass[10pt,a4paper]{article}

\usepackage[T1]{fontenc}                          % FONT
\usepackage{lmodern}                              % FONT
\renewcommand{\familydefault}{\ttdefault}         % FONT: whole document in LM Mono

\usepackage[margin=2cm]{geometry}
\usepackage{xcolor}
\usepackage{longtable}
\usepackage{booktabs}
\usepackage{array}
\usepackage{fancyhdr}
\usepackage[hidelinks]{hyperref}

% Monospace sets poorly when justified; ragged-right keeps word spacing constant.
\raggedright
\setlength{\parindent}{0pt}
\setlength{\parskip}{0.55em}
\newcommand{\tightlist}{\setlength{\itemsep}{1pt}\setlength{\parskip}{0pt}\setlength{\parsep}{0pt}}

\definecolor{rule}{gray}{0.55}
\definecolor{box}{gray}{0.94}
\newcommand{\us}{\_}
\newcommand{\cmd}[1]{\colorbox{box}{#1}}
\newcommand{\done}{[x]}
\newcommand{\todo}{[\phantom{x}]}
\newcommand{\Sh}{\textbf{S}}
\newcommand{\Mu}{\textbf{M}}

% Sections are the four Parts; number them that way rather than 1-4 plus a title.
\renewcommand{\thesection}{Part \arabic{section}}
\renewcommand{\thesubsection}{\arabic{section}.\arabic{subsection}}
\makeatletter
\renewcommand*\l@section{\@dottedtocline{1}{0em}{5.5em}}
\renewcommand*\l@subsection{\@dottedtocline{2}{2em}{3em}}
\makeatother

\pagestyle{fancy}
\fancyhf{}
\lhead{MOSIP Decode 2026 / PS1 / Project Handbook}
\rhead{\thepage}
\renewcommand{\headrulewidth}{0.4pt}

\begin{document}

\begin{titlepage}
\vspace*{3cm}
{\Large\bfseries MOSIP Decode 2026 --- Problem Statement 1}\\[0.6em]
{\large Automated Conformance Testing for Inji Certify and Inji Verify}\\[2em]
{\LARGE\bfseries Project Handbook}\\[2em]
Shardul Chogale (conformance engine and evidence)\\
Mukta Varak (platform and delivery)\\[2em]
Written: 30 September 2026; updated 2 October 2026\\
Final submission deadline: \textbf{26 October 2026, 11:59 PM IST}\\[3em]
\begin{tabular}{@{}ll@{}}
Part 1 & Mentor session brief --- Weekly Connect, 7 October\\
Part 2 & Start to end: the clean guide to running the project\\
Part 3 & Record of every step taken so far\\
Part 4 & Planner, 30 September to 26 October, for both of us\\
\end{tabular}
\vfill
{\small This is the document to open first. It is kept current: when a procedure changes,
Part 2 changes with it; when a day's work is done, Part 3 gains an entry and Part 4 gets
its boxes ticked. The detailed evidence behind every claim lives in
docs/findings/findings.tex.}
\end{titlepage}

\tableofcontents
\newpage

%==========================================================================
\section{Mentor session brief}
%==========================================================================

\begin{tabular}{@{}ll@{}}
Session & MOSIP Decode Weekly Connect\\
When & Wednesday 7 October 2026, 5:00--6:00 PM IST\\
Where & \url{https://mosip.zoom.us/j/89100675963}\\
Attending & Shardul (and Mukta, if available)\\
\end{tabular}

\subsection{Where we are, in one paragraph}

We have a working harness that drives the OpenID Foundation conformance suite's REST API
against Inji Verify with no human interaction. One command ---
\cmd{./run-conformance.sh --component verify} --- creates the test plan, starts each
module, obtains an authorisation request from Inji Verify's own API, delivers it to the
suite, collects the result, and gates it against a recorded baseline. On 29 September its
first unattended run reproduced our Day 1 manual run \textbf{check for check}: 15
passing checks, 4 failures, 3 warnings, the identical set of seven non-passing checks, and
the same point of termination. Versions tested: Inji Verify 0.18.2, conformance suite
5.3.1. Inji Certify is not yet running.

\textbf{Since then (2 October):} we found that upstream has an Inji \textbf{1.0 line}
(\texttt{release-1.0.x}, published as \texttt{1.0.0-alpha.1}) --- which is what the problem
statement means by ``on 1.0''. On it, DCQL is implemented and the nonce is random, so our
three headline findings do not apply there. We are moving the baseline to
\texttt{1.0.0-alpha.1}; the findings stand for 0.18.2 and are on hold for upstream
reporting.

\subsection{Three things to show if asked}

\begin{enumerate}\tightlist
  \item \textbf{The automated run matches the manual run exactly.} This is what makes
        everything else credible: the harness measures what a tester would see by hand.
  \item \textbf{The nonce experiment.} Two runs differing only in where the nonce comes
        from. With the nonce the Verify UI's SDK sends, the entropy check warns (71.68
        bits against 96 required) and the URL-safety check fails. When verify-service
        generates the nonce itself, both pass and nothing else changes. The weak nonce is
        \texttt{btoa(Date.now())} in the SDK; the service already does it correctly.
  \item \textbf{The benchmark gate.} Inji Verify fails every module of the OID4VP 1.0
        Final plan at DCQL extraction, so the gate fails a build on \emph{change from the
        recorded baseline}, not on absolute pass rate. The suite's own verdict is never
        rewritten --- only the TestNG verdict is gated.
\end{enumerate}

\subsection{Questions, ranked by how much the answer changes the build}

Each question has a default we will follow if there is no answer, so nothing blocks on
the session. The top three are the ones worth the time.

\begin{longtable}{@{}p{0.5cm}p{5.0cm}p{5.0cm}p{5.0cm}@{}}
\toprule
\# & Question & Why it matters & Our default if unanswered\\
\midrule
\endhead
1 & \textbf{Q1 --- Which spec version is the benchmark?} Verify 0.18.2 implements
Presentation Exchange; OID4VP 1.0 Final requires DCQL.
& Decides whether any module can pass, and what the gate measures. Every 1.0 Final module
currently stops at DCQL extraction.
& 1.0 Final. Since the 1.0 line implements DCQL, this is now the natural benchmark rather
than an all-fail one; the ID2 plan drops to optional.\\
\addlinespace
2 & \textbf{Confirm the 1.0 target.} We found \texttt{release-1.0.x} (Verify and Certify) and the
published \texttt{1.0.0-alpha.1}; \texttt{alpha.2} is named in the branch but not yet published.
Is \texttt{1.0.0-alpha.1} the intended target, or should we build the branch from source? Is
the 0.18.x line still maintained?
& Decides what our baseline is, and whether the 0.18.2 findings are worth reporting at all.
& Target \texttt{1.0.0-alpha.1}; hold the 0.18.2 findings unless told the line is
maintained.\\
\addlinespace
3 & \textbf{Q7 --- Is there a lighter Inji Certify setup?} \texttt{docker-compose-injistack}
needs a PKCS12 keystore from Mimoto onboarding, a public DID endpoint, an external
authorisation server and partner API keys.
& Decides whether the combined run is achievable locally. Certify's issuer plan is 61
modules against the verifier plan's 12.
& Attempt the local stack; if blocked, submit Verify end to end plus a documented Certify
path, and say so plainly.\\
\addlinespace
4 & \textbf{Where should findings go upstream?} GitHub issues on
\texttt{mosip/inji-verify}, or another tracker? Would a pull request fixing the SDK nonce
be welcome from us?
& This is our open-source contribution path. Our strongest finding is a three-line fix.
& A GitHub issue on \texttt{mosip/inji-verify}, followed by a pull request from our
fork.\\
\addlinespace
5 & \textbf{Q4 --- Benchmark semantics.} Regression-from-baseline, or an absolute target the
build fails until conformance improves?
& Changes what ``pass'' means in CI.
& Regression-from-baseline (built). Absolute is one flag: \texttt{--policy
absolute-pass-rate}.\\
\addlinespace
6 & \textbf{Submission criteria.} The kickoff settled the list --- source code, working
prototype, slide deck, video demo. Remaining: must the api-testrig classes be merged upstream,
or are branches on our forks acceptable?
& Scopes the last week.
& Deliver the Java classes as fork branches shaped as upstream pull requests.\\
\addlinespace
7 & \textbf{Q2 --- HAIP.} Verify hard-codes \texttt{direct\us{}post}; HAIP requires
\texttt{direct\us{}post.jwt}. Track HAIP as a known-failing suite, or exclude it?
& Affects scope, not architecture.
& Exclude from the gate; document the gap.\\
\addlinespace
8 & \textbf{Q3 --- Client identifier schemes.} Verify offers \texttt{did} and
\texttt{pre\us{}registered}; the 1.0 Final plan offers \texttt{redirect\us{}uri},
\texttt{x509\us{}san\us{}dns}, \texttt{x509\us{}hash}. No overlap.
& Explains one of the failures in every run.
& Record as a known gap.\\
\addlinespace
9 & \textbf{Q5 / Q8 --- Configuration questions.} Is the shipped \texttt{config.json}
meant to represent a certifiable deployment? Does anything rely on the presentation definition id,
which the UI hard-codes?
& Decides whether two observations become findings.
& Keep both as observations, not findings.\\
\addlinespace
--- & \textbf{Q6 --- answered by us; confirm only.} Unattended runs work: the harness
delivers the request by HTTP GET to the test's exposed \texttt{authorization\us{}endpoint},
which is what the suite's paste box does internally.
& Confirms the approach is acceptable.
& Proceed.\\
\bottomrule
\end{longtable}

\subsection{What we want to leave with}

At minimum: confirmation of the 1.0 target (question 2), the Certify route (question 3),
and whether the 0.18.2 findings are worth reporting and where (question 4). Everything else
has a workable default.

\subsection{Notes during the session}

\vspace{0.3em}
\begin{tabular}{@{}p{1.4cm}p{14.6cm}@{}}
Q1 & \hrulefill\\[1.1em]
Ver. & \hrulefill\\[1.1em]
Q7 & \hrulefill\\[1.1em]
Issues & \hrulefill\\[1.1em]
Other & \hrulefill\\[1.1em]
 & \hrulefill\\
\end{tabular}

After the session, record the answers in \texttt{docs/findings/findings.tex} (the Open
Questions section) and update the planner in Part 4 where an answer changes it.

\newpage
%==========================================================================
\section{Start to end: running the project from a clean machine}
%==========================================================================

Follow this top to bottom on a machine that has never run the project. Every command has
been run on the reference machine (macOS, Apple Silicon); Windows differences are noted
where they exist. Budget two to three hours the first time, most of it image downloads.

\subsection{What runs, and why each piece is there}

\begin{longtable}{@{}p{3.4cm}p{5.8cm}p{6.6cm}@{}}
\toprule
Piece & Where it listens & Why it is needed\\
\midrule
\endhead
OpenID conformance suite 5.3.1 & \url{https://localhost.emobix.co.uk:8443} & The examiner. It plays
the wallet (verifier tests) or the client (issuer tests) and checks the component against
the spec.\\
Inji Verify 1.0.0-alpha.1 & UI :3000, API :8080 (path /v1/verify) & The verifier under test. Happy-flow passes on it as of 2 October.\\
Inji Certify 0.14.0 & not yet running & The issuer under test (Part 4, weeks 3--4).\\
ngrok tunnel & a *.ngrok-free.dev address & Mandatory. Verify identifies itself with a
did:web DID resolved over public HTTPS, and bakes its public address into the response\us{}uri
the suite must reach from inside its own container.\\
The runner (this repo) & a command, not a server & Drives the suite's REST API and
Verify's API, normalises results, applies the gate.\\
\bottomrule
\end{longtable}

\subsection{Step 1 --- Prerequisites}

Git, Docker Desktop, Java 21, Maven 3.9+, Python 3.11+, ngrok. Check:

\begin{verbatim}
git --version; docker --version; java -version; mvn -version
python3 --version; ngrok version
\end{verbatim}

macOS: \cmd{brew install git openjdk@21 maven python3 ngrok}, plus Docker Desktop from
docker.com. Windows: winget for the same set, Docker Desktop with WSL 2, and once only:
\cmd{git config --global core.autocrlf input} so shell scripts keep Unix line endings.

Java is \textbf{21}, not 11: the Inji api-test projects require it, whatever older MOSIP
documentation says.

\subsection{Step 2 --- The working directory}

The runner expects these six folders side by side under one root. The folder names
matter: the runner finds Inji Verify's UI configuration at \texttt{../inji-verify}.

\begin{verbatim}
mkdir -p ~/projects/MOSIP && cd ~/projects/MOSIP
git clone https://github.com/shard-c6/mosip-decode-ps1.git
git clone --depth 1 https://github.com/mosip/inji-verify.git
git clone --depth 1 https://github.com/mosip/inji-certify.git
git clone --depth 1 https://github.com/mosip/mosip-functional-tests.git
git clone https://gitlab.com/openid/conformance-suite.git
git clone https://gitlab.com/openid/conformance-suite-automated-testing-tutorial.git
\end{verbatim}

Then the \textbf{target version}, Inji Verify 1.0, as a seventh folder: a worktree of the
\texttt{inji-verify} clone at the published 1.0 tag. The 0.18.2 checkout stays as it is,
because it reproduces the earlier evidence.

\begin{verbatim}
git -C inji-verify fetch --depth 1 https://github.com/mosip/inji-verify.git \
    tag v1.0.0-alpha.1
git -C inji-verify worktree add --detach ../inji-verify-1.0 v1.0.0-alpha.1
\end{verbatim}

\texttt{mosip-decode-ps1} is ours. The Java work lands on branches of our forks of
\texttt{inji-verify} and \texttt{inji-certify} (Shardul's forks at
github.com/shard-c6; Mukta has write access to both).

\subsection{Step 3 --- ngrok, once}

Create a free ngrok account, then \cmd{ngrok config add-authtoken <token>} and claim the
account's free \textbf{static} domain in the dashboard. Without a static domain the
hostname changes on every restart and Step 4 must be redone each time. The domain is
per-developer: Shardul's is sphere-dust-sandlot.ngrok-free.dev; Mukta needs her own.

\subsection{Step 4 --- Three local edits to Inji Verify 1.0}

In \texttt{inji-verify-1.0/docker-compose/}. None is committed anywhere; each one missing
produces a failure that looks like a defect in Inji Verify. Each was established by the run
that needed it (findings, Day 6).

\textbf{4a. Public host.} The literal \texttt{VERIFY\us{}SERVICE\us{}PROXY\us{}FOR\us{}LOCALHOST}
appears in five values of \texttt{docker-compose.yml} and is substituted by nothing. Replace it
with your domain --- hostname only:

\begin{verbatim}
sed -i '' 's/VERIFY_SERVICE_PROXY_FOR_LOCALHOST/YOUR-DOMAIN.ngrok-free.dev/g' \
    docker-compose.yml                  # macOS; on Linux drop the ''
grep -c "YOUR-DOMAIN" docker-compose.yml      # expect 5
\end{verbatim}

Windows (PowerShell):
\begin{verbatim}
(Get-Content docker-compose.yml) -replace 'VERIFY_SERVICE_PROXY_FOR_LOCALHOST',
  'YOUR-DOMAIN.ngrok-free.dev' | Set-Content docker-compose.yml
\end{verbatim}

\textbf{4b. Database script (finding F-10).} In \texttt{db-init/init.sql}, replace the
\texttt{verify.vp\us{}submission} table with the definition in
\texttt{../db\us{}scripts/inji\us{}verify/ddl/verify-vp\us{}submission.sql} (keep the
\texttt{verify.} schema prefix). Without this, every presentation submission returns HTTP 500.
If the stack has already run once, the old schema is in the data volume: recreate it with
\cmd{docker compose up -d --force-recreate --renew-anon-volumes postgres}.

\textbf{4c. A conformance credential.} In \texttt{config/config.json}, add to
\texttt{verifiableClaims}:

\begin{verbatim}
{ "logo": "/assets/cert.png",
  "name": "OIDF Conformance PID (SD JWT)",
  "clientIdPrefix": "pre_registered",
  "purpose": "OpenID Foundation conformance testing",
  "dcqlQuery": { "credentials": [ {
      "id": "pid_sd_jwt", "format": "dc+sd-jwt",
      "meta": { "vct_values": [ "urn:eudi:pid:1" ] } } ] } }
\end{verbatim}

The suite's verifier tests always present a credential of type \texttt{urn:eudi:pid:1}; a
verifier that asks for anything else rejects it, correctly. \texttt{pre\us{}registered} makes
Verify deliver the request inline, which the test variant needs.

\subsection{Step 5 --- Point the runner at your machine}

\textbf{5a. A signing key for the suite's test credential.} Private, so it lives in the
gitignored \texttt{configs/keys/}. Generated with the suite's own script:

\begin{verbatim}
cd ~/projects/MOSIP/mosip-decode-ps1
python3 ../conformance-suite/scripts/generate-vp-test-cert.py \
    --hostname YOUR-DOMAIN.ngrok-free.dev \
    --output configs/keys/vp-signing-jwk.json --ca-output configs/keys/vp-test-ca.pem
\end{verbatim}

\textbf{5b. Your own settings.} \texttt{configs/runner.json} is shared and committed. Anything
that is yours goes in \texttt{configs/runner.local.json}, which git ignores and the runner
deep-merges on top, matching components by name. The client identifier must use the
\texttt{redirect\us{}uri:} prefix with \emph{your} response URI, and a distinct alias keeps two
developers' plans apart:

\begin{verbatim}
{ "components": [ { "component": "inji-verify", "alias": "injiverify-YOURNAME",
    "verifierRequest": { "clientId":
      "redirect_uri:https://HOST/v1/verify/v2/vp-submission/direct-post"
} } ] }
\end{verbatim}

\texttt{HOST} is your ngrok domain, e.g. \texttt{sphere-dust-sandlot.ngrok-free.dev}.

Install the runner's one dependency: \cmd{python3 -m pip install -r requirements.txt}.

\subsection{Step 6 --- Start the stack, in this order}

\textbf{6a. Docker Desktop.} Start it and wait for the engine: \cmd{docker info}.

\textbf{6b. Conformance suite.} Note the \texttt{-f}: without it, stop and start target
the wrong file and report success while doing nothing.

\begin{verbatim}
cd ~/projects/MOSIP/conformance-suite
docker compose -f docker-compose-prebuilt.yml up -d
\end{verbatim}

\textbf{6c. Inji Verify 1.0.} Only one Verify version can run at a time --- they share
container names and ports --- so take the other one down first.

\begin{verbatim}
cd ~/projects/MOSIP/inji-verify/docker-compose && docker compose down     # 0.18.2, if up
cd ~/projects/MOSIP/inji-verify-1.0/docker-compose && docker compose up -d
\end{verbatim}

\textbf{6d. The tunnel} --- leave it running in its own terminal:
\begin{verbatim}
ngrok http --url=https://YOUR-DOMAIN.ngrok-free.dev 3000
\end{verbatim}

\subsection{Step 7 --- Check everything is healthy}

All three must answer before a run is worth attempting:

\begin{verbatim}
curl -k https://localhost.emobix.co.uk:8443/api/currentuser    # JSON, dev user
curl http://localhost:8080/v1/verify/did.json                  # Verify, local
curl https://YOUR-DOMAIN.ngrok-free.dev/v1/verify/did.json     # Verify, public
\end{verbatim}

The last is the one that matters most: if the public DID document does not resolve,
nothing downstream can work. The path is /v1/verify/did.json --- there is no
.well-known, because the DID has path segments.

\subsection{Step 8 --- Run the unit tests}

Needs no Docker. Normalisation is tested against the real Day 1 log, so these also check
that the suite's field names have not drifted from our contract.

\begin{verbatim}
cd ~/projects/MOSIP/mosip-decode-ps1 && python3 -m pytest tests/ -q     # expect 63 passed
\end{verbatim}

\subsection{Step 9 --- Run conformance}

\begin{verbatim}
./run-conformance.sh --component verify                  # the full verifier plan
./run-conformance.sh --component verify \
    --only oid4vp-1final-verifier-happy-flow -v          # one module, verbose
./run-conformance.sh --component verify --output scratch/run.json
./run-conformance.sh --component verify \
    --config configs/runner-0.18.json                    # 0.18.2: earlier evidence
./run-conformance.sh --component certify                 # once Certify is configured
./run-conformance.sh --combined                          # every configured component
\end{verbatim}

Useful options: \texttt{--nonce-mode sdk|service} (Part 3, 29 September, for why this
exists), \texttt{--policy absolute-pass-rate}, \texttt{--no-gate-exit} for exploratory
runs only.

\textbf{Exit status is the verdict}, and CI keys on it:

\begin{tabular}{@{}ll@{}}
0 & gate passed --- nothing regressed against the baseline\\
1 & gate failed --- a regression, or the harness could not execute a module\\
2 & configuration error --- nothing was run\\
\end{tabular}

\subsection{Step 10 --- Verify the results}

A run that finished is not yet a run that is right. Check, in order:

\begin{enumerate}\tightlist
  \item \textbf{The gate summary} printed at the end. ``Harness could not execute''
        means a fault in us or the environment, never a conformance result.
  \item \textbf{Each module's handoff.} In the output JSON, verifier modules carry a
        \texttt{handoff} block: \texttt{delivered} true, \texttt{httpStatus} 302, and
        \texttt{secondsRemaining} well above zero. The request expires 300 s after issue.
  \item \textbf{Checks against the baseline.} Every non-passing check carries a
        \texttt{findingRef} (F-01 \ldots{} F-09) when it matches a known finding. A failing
        check with no \texttt{findingRef} is new and needs reading before anything else.
  \item \textbf{The suite's own view} at the \texttt{logUrl} for any module in doubt.
  \item \textbf{Archived logs} under \texttt{logs/<date>/}. These are the JSON the runner
        fetched and are \emph{unsigned}. For anything that will be reported upstream,
        export the signed log from the suite (\emph{Download all Logs}) and commit the
        \texttt{.json} and \texttt{.sig} together.
\end{enumerate}

\subsection{Step 11 --- Record what the run showed}

\begin{itemize}\tightlist
  \item A new or changed finding: a \texttt{\textbackslash finding} block in
        \texttt{docs/findings/findings.tex}, tagged CONFIRMED or UNVERIFIED, with the log
        line quoted and the spec section the suite cited. Nothing goes upstream while
        UNVERIFIED.
  \item A day's work: a dated subsection in the findings Day Log --- objective, runs
        with test IDs, decisions with reasons, status.
  \item A module whose result is now understood: an entry in
        \texttt{configs/contract/expected-failures.json} with its \texttt{reason} and an
        \texttt{evidence} path. An expected failure without a reason is a silenced test.
  \item Every run worth keeping: a row in \texttt{logs/README.md}.
\end{itemize}

Rebuild the findings PDF: \cmd{cd docs/findings \&\& pdflatex findings.tex \&\& pdflatex
findings.tex}.

\subsection{Step 12 --- Commit}

Stage by path, never blindly: \cmd{git status --short --untracked-files=all} first.
Subject in the imperative; body explaining why. \textbf{No assistant attribution lines of
any kind} --- commits are made under our own names. Push to \texttt{master} on
\texttt{shard-c6/mosip-decode-ps1}.

\subsection{Step 13 --- Shut down}

Leaving the stack running costs memory and is the commonest cause of ``it worked
yesterday''.

\begin{verbatim}
cd ~/projects/MOSIP/inji-verify/docker-compose && docker compose down
cd ~/projects/MOSIP/conformance-suite && docker compose -f docker-compose-prebuilt.yml down
docker ps                        # expect headers only
\end{verbatim}

Then stop ngrok (Ctrl-C in its terminal). \texttt{down} keeps volumes, so the suite's
plans and logs survive. \textbf{Never} add \texttt{-v} to the suite's \texttt{down} unless
you mean to delete every plan and log it holds.

\subsection{Step 14 --- Ending the project}

The final week, in order. Detail and dates in Part 4.

\begin{enumerate}\tightlist
  \item Fresh-clone test: follow only this Part 2 on a machine that has never run the
        project. Anything missing gets fixed here, not explained later.
  \item Final combined run, recorded; signed logs exported and committed.
  \item Upstream: file confirmed findings as issues against \texttt{mosip/inji-verify},
        and offer the fixes as pull requests from our forks.
  \item Submission package: public repo with a clear README; the 1--2 page document or
        5--7 slide pitch; demo video; the Docker Compose, test plan configuration,
        one-command runner and GitHub Actions workflow the problem statement lists.
  \item Submit early (target 23 October). Multiple submissions are allowed and only the
        last before 26 October, 11:59 PM IST counts, so an early one is a safety net.
  \item Hand-over: MOSIP may request the code be contributed. Keep the repository
        clean enough that this is a formality.
  \item Shut the stack down (Step 13).
\end{enumerate}

\subsection{Troubleshooting --- problems that have already cost us time}

\begin{longtable}{@{}p{5.4cm}p{11cm}@{}}
\toprule
Symptom & Cause and fix\\
\midrule
\endhead
Verify returns HTTP 400 to vp-session-request & The presentation definition has no
\texttt{id}. The runner builds the UI's envelope; if this appears, check the credential
name in \texttt{runner.json} matches \texttt{config.json} exactly.\\
A module fails at request\us{}uri & The credential uses \texttt{did} (by reference). Redo
Step 4b.\\
Failure that looks like an expired request & Requests die 300 s after issue. The runner
refuses to deliver with under 30 s left; manually, check liveness with
\texttt{curl .../v1/verify/vp-request/<id>} first.\\
Verify behaves as if unconfigured & Environment variables are read at container
\emph{creation}. After editing docker-compose.yml: \texttt{down}, then \texttt{up -d}.
\texttt{restart} is not enough.\\
Suite ``stopped'' but still running & A plain \texttt{docker compose down} in that folder
targets the wrong file. Always \texttt{-f docker-compose-prebuilt.yml}.\\
DID document 404 & The path is /v1/verify/did.json; no .well-known.\\
Gate passed but nothing ran & Fixed on 29 September: a harness error now fails the gate.
If you ever see this again, it is a regression in the gate.\\
Unit tests pass, CLI crashes on import & Fixed on 29 September by import tests. Run the
full \texttt{pytest tests/}, not a subset.\\
Browser warns the suite is not secure & Expected: self-signed certificate on a hostname
that resolves to 127.0.0.1. Applies only because it is your own machine.\\
\bottomrule
\end{longtable}

\newpage
%==========================================================================
\section{Record of every step taken}
%==========================================================================

What we did, in order, and what each step produced. Detail, evidence and quoted logs for
every item are in \texttt{docs/findings/findings.tex}; this is the index to it.

\subsection{16--20 September --- registration and pre-kickoff groundwork}

\begin{enumerate}\tightlist
  \item Registered as a team for MOSIP Decode 2026 (confirmed 17 September) and chose
        Problem Statement 1.
  \item Cloned the five upstream repositories side by side under ~/projects/MOSIP.
  \item Started the conformance suite 5.3.1 locally from prebuilt images. Established
        that its local profile needs no API token.
  \item Started Inji Verify 0.18.2 (three containers). Found the unsubstituted
        \texttt{VERIFY\us{}SERVICE\us{}PROXY\us{}FOR\us{}LOCALHOST} placeholder and replaced
        it with an ngrok hostname; found that did:web makes a public tunnel mandatory.
  \item Created the verifier test plan \texttt{lN7C4DH1HvYmc} (OID4VP 1.0 Final alpha,
        sd\us{}jwt\us{}vc / redirect\us{}uri / url\us{}query / plain\us{}vp / direct\us{}post).
  \item Ran the happy-flow module by hand three times. Two runs were discarded, each for
        a reason we could prove: an expired request, then a variant mismatch fixed by
        switching the SD-JWT credential to \texttt{pre\us{}registered}. Both had looked like
        Inji Verify defects. The third, \texttt{MtzpWoD2dVgtic0}, is the baseline.
  \item Catalogued nine findings (F-01 to F-09), eight confirmed. Headline: Verify
        implements Presentation Exchange, not DCQL, so no 1.0 Final module can pass.
  \item Wrote \texttt{build-oid4vp-uri.py}, the first piece of tooling; scaffolded the
        repository with setup, runbook, plan and findings documents.
\end{enumerate}

\subsection{23 September --- official kickoff: plan, roles, contract}

\begin{enumerate}\tightlist
  \item Re-anchored the plan: the 20 September work predated the event.
  \item Agreed roles by skill rather than by language. Mukta: platform and delivery.
        Shardul: conformance engine and evidence. Estimated split 49/51.
  \item Chose to put the api-testrig work on our forks of inji-verify and inji-certify,
        shaped as upstream pull requests. api-test lives in those repositories, not in
        mosip-functional-tests.
  \item Wrote the runner-to-testrig contract (\texttt{docs/CONTRACT.md}) with a fixture
        generated from the real Day 1 log, so the Java side need not wait for the runner.
  \item Traced F-01/F-02 to \texttt{btoa(Date.now())} in \texttt{inji-verify-sdk}, and
        found that verify-service already generates a sound nonce but only as a fallback.
        Drafted the upstream issue.
\end{enumerate}

\subsection{24 September --- the runner}

\begin{enumerate}\tightlist
  \item Resolved the week's blocking unknown from the suite's source: a verifier module
        receives its authorisation request by an ordinary HTTP GET to the
        \texttt{authorization\us{}endpoint} it exposes. No browser is needed.
  \item Found the suite's built-in path for the screenshot requirement: a browser
        automation entry against an HTML stand-in page the suite serves itself.
  \item Built the runner: configuration, suite client, Verify client, handoff,
        normalisation, gate, orchestrator, CLI, and \texttt{run-conformance.sh}.
  \item 24 unit tests; one caught a defect in the gate before any live run.
\end{enumerate}

\subsection{29 September --- first live runs}

\begin{enumerate}\tightlist
  \item Found and fixed two request defects before running, by reading Verify's DTO:
        no presentation definition was sent, and no nonce --- which would have made F-01
        and F-02 silently vanish from automated runs.
  \item Run 1 failed with HTTP 400: the definition lacked an \texttt{id}. Traced to the
        envelope the UI builds in \texttt{vpVerificationState.ts}; the runner now builds
        the same one.
  \item Run 1 also exposed the most serious defect: the gate reported PASSED having
        measured nothing. Harness errors now always fail the gate.
  \item Run 2, \texttt{hcuWEvxGbvRgt7e}: first fully unattended run. Matches the manual
        baseline check for check.
  \item Run 3, \texttt{Q9MaDZyfpJhjiDp}: controlled experiment on the nonce source.
        Confirmed the F-01/F-02 root cause.
  \item Emailed Mukta the directory setup and project status. 50 unit tests.
\end{enumerate}

\subsection{30 September --- deadline found, this handbook}

\begin{enumerate}\tightlist
  \item Found the real deadline on the official event page: \textbf{26 October 2026,
        11:59 PM IST}, not the 21 October we had assumed. The plan gains five days.
  \item Found the Weekly Connect schedule --- Wednesdays, 5:00--6:00 PM IST, from today.
  \item Wrote this handbook: mentor brief, start-to-end guide, this record, the planner.
\end{enumerate}

\subsection{2 October --- the kickoff recording, and the 1.0 line}

\begin{enumerate}\tightlist
  \item Worked through the transcript of the 23 September kickoff webinar, which we had not
        attended live. Correction: earlier documents said it had not taken place; it had.
  \item Took from it: the mandatory submissions are source code, a working prototype, a
        slide deck and a video demo; the Docker Compose must bring up each module alone or
        both together; CI's purpose is catching non-conformance per change.
  \item Followed up a passing mention of ``DCQL requests'' and found the Inji 1.0 line:
        \texttt{release-1.0.x} in Verify and Certify, published as \texttt{1.0.0-alpha.1}.
        On it, DCQL exists and the nonce is random --- F-01, F-02 and F-03 do not apply.
  \item Put the F-01/F-02 upstream issue on hold; scoped every finding to 0.18.2; decided to
        move the baseline to \texttt{1.0.0-alpha.1}.
  \item Taught the runner the 1.0 API (DCQL, \texttt{/v2/vp-session-request}, the 1.0 SDK's
        behaviour). Resolved F-05 as our own artefact on the way.
  \item Five runs on \texttt{1.0.0-alpha.1}, each changing one thing. Found F-10 (the Compose
        database script breaks every submission) and F-11 (sample config uses a superseded
        format). Established what a passing run needs.
  \item \textbf{Happy-flow passed: FINISHED / REVIEW, 35 checks, 0 failures, 0 warnings, fully
        unattended} --- including the screenshot. The gate now protects a real pass.
  \item Ran the full plan. Added REVIEW resolution: the harness fetches Inji Verify's own
        verdict, because the suite cannot see whether a forged credential was rejected.
        Result: 9 of 11 modules exercised, all consistent with conformance.
\end{enumerate}

\newpage
%==========================================================================
\section{Planner, 30 September to 26 October}
%==========================================================================

\Sh{} = Shardul (conformance engine and evidence). \Mu{} = Mukta (platform and delivery).
Tick boxes as work lands; move a slipped item forward rather than deleting it, so the
planner stays an honest record.

\subsection{Fixed points}

\begin{tabular}{@{}p{3.4cm}p{13cm}@{}}
Wed 30 Sep & Weekly Connect (first session)\\
Sun 4 Oct & End of week 2: verifier plan baselined on \texttt{1.0.0-alpha.1}; api-test builds\\
Wed 7 Oct & Weekly Connect --- mentor brief, Part 1\\
Sun 11 Oct & End of week 3: runner output lands in the api-testrig report\\
Wed 14 Oct & Weekly Connect\\
Sun 18 Oct & End of week 4: \textbf{minimum viable submission}\\
Tue 20 Oct & Nothing new or risky starts after today\\
Wed 21 Oct & Weekly Connect --- show the working build if possible\\
Fri 23 Oct & \textbf{First submission} (safety net)\\
Mon 26 Oct & \textbf{Final submission, 11:59 PM IST} --- aim to finish by 6:00 PM\\
\end{tabular}

\subsection{Week 2 remainder --- Wed 30 Sep to Sun 4 Oct: move to the 1.0 line}

Re-planned on 2 October. Wednesday and Thursday's items did not happen and are carried
forward; the discovery of the 1.0 line changes what ``finish the verifier side'' means.

\begin{longtable}{@{}p{1.9cm}p{6.9cm}p{7.2cm}@{}}
\toprule
Day & \Sh & \Mu\\
\midrule
\endhead
Wed 30 & \todo{} Weekly Connect (attendance not recorded)
& \done{} Directory setup email received (29 Sep)\newline
\done{} Sent the kickoff recording link (30 Sep)\\
\addlinespace
Fri 2 Oct & \done{} Kickoff transcript worked through\newline
\done{} Found the 1.0 line; findings scoped; issue on hold\newline
\done{} Bring up Verify \texttt{1.0.0-alpha.1}; run happy-flow --- \textbf{passes}\newline
\done{} F-10 found; upstream issue drafted
& \todo{} Accept the two fork invitations; claim an ngrok domain\newline
\todo{} Build apitest-commons; api-test in inji-verify\newline
\textbf{Build api-test from \texttt{release-1.0.x}}, not master --- the target changed\\
\addlinespace
Sat 3 Oct & \done{} Full 1.0 Final plan on \texttt{alpha.1} --- done early, 2 Oct: 9 of 11
exercised and consistent, 2 n/a to the variant; baseline for all 11\newline
\todo{} Make \texttt{nonceMode} follow the version under test
& \todo{} api-test in inji-certify (\texttt{release-1.0.x})\newline
\todo{} Write down the TestNG/Extent result shape\\
\addlinespace
Sun 4 Oct & \done{} Screenshot mechanism --- confirmed by run, 2 Oct\newline
\todo{} Re-check F-05, F-06, F-07, F-09 on 1.0
& \todo{} Review CONTRACT.md; OpenIDConformanceTest skeleton\newline
\todo{} Read the 1.0 compose; sketch per-module profiles\\
\bottomrule
\end{longtable}

\textbf{Exit check:} \texttt{./run-conformance.sh --component verify} runs all 12 modules
against \texttt{1.0.0-alpha.1} with a baseline entry for each; api-test builds in both
repositories on \texttt{release-1.0.x}.

\subsection{Week 3 --- Mon 5 to Sun 11 Oct: the two halves meet}

\begin{longtable}{@{}p{1.9cm}p{6.9cm}p{7.2cm}@{}}
\toprule
Day & \Sh & \Mu\\
\midrule
\endhead
Mon 5 & \todo{} Issuer-role support in the runner: no handoff; different plan
configuration & \todo{} Unified Docker Compose with \textbf{per-module profiles}: Verify, Certify,
or both, each with the suite --- the kickoff's stated requirement\\
\addlinespace
Tue 6 & \todo{} Issuer plan configuration templates\newline
\todo{} Upstream reporting decision, per the mentor answer on 0.18.x
& \todo{} Compose: solve cross-project networking and the public hostname\\
\addlinespace
Wed 7 & \todo{} Weekly Connect --- brief in Part 1\newline
\todo{} File any finding that survives on 1.0
& \todo{} OpenIDConformanceTest in inji-verify's api-test, reading real runner
output\\
\addlinespace
Thu 8 & \todo{} Combined mode: consolidated output across components
& \todo{} Inji Certify environment: keystore, DID endpoint, auth server, keys --- day 1\\
\addlinespace
Fri 9 & \todo{} Signed exports; findings catch-up
& \todo{} Certify environment --- day 2\\
\addlinespace
Sat 10 & \todo{} Architecture and automation documentation
& \todo{} Certify environment --- day 3, or fall back (see risks)\\
\addlinespace
Sun 11 & \multicolumn{2}{p{14.2cm}}{\todo{} Joint integration session: runner output into the Extent
report, end to end. Buffer otherwise.}\\
\bottomrule
\end{longtable}

\textbf{Exit check:} a Verify run's results appear as TestNG pass/fail/skip in Inji
Verify's own api-testrig report.

\subsection{Week 4 --- Mon 12 to Sun 18 Oct: complete the mandatory set}

\begin{longtable}{@{}p{1.9cm}p{6.9cm}p{7.2cm}@{}}
\toprule
Day & \Sh & \Mu\\
\midrule
\endhead
Mon 12 & \todo{} Run the issuer plan against Certify & \todo{} Fold Certify into the unified
compose\\
\addlinespace
Tue 13 & \todo{} Triage the issuer results; baseline for Certify
& \todo{} OpenIDConformanceTest in inji-certify's api-test\\
\addlinespace
Wed 14 & \todo{} Weekly Connect\newline \todo{} Combined run end to end
& \todo{} GitHub Actions workflow, first version\\
\addlinespace
Thu 15 & \todo{} Self-certification documentation
& \todo{} Result diff between runs; consolidated report and badge\\
\addlinespace
Fri 16 & \todo{} Selective execution polish; expected-failures handling
& \todo{} Actions: fail the build on a regression, demonstrably\\
\addlinespace
Sat 17 & \todo{} Findings written up as issue drafts
& \todo{} Setup, troubleshooting and CI documentation\\
\addlinespace
Sun 18 & \multicolumn{2}{p{14.2cm}}{\todo{} \textbf{Minimum viable submission check}: all three
invocations work and land in the api-testrig output.}\\
\bottomrule
\end{longtable}

\subsection{Week 5 --- Mon 19 to Mon 26 Oct: submit}

\begin{longtable}{@{}p{1.9cm}p{6.9cm}p{7.2cm}@{}}
\toprule
Day & \Sh & \Mu\\
\midrule
\endhead
Mon 19 & \todo{} Documentation complete; README final
& \todo{} Fresh-clone test using only Part 2\\
\addlinespace
Tue 20 & \todo{} Fix whatever the fresh-clone test found
& \todo{} Fix whatever the fresh-clone test found\\
\addlinespace
Wed 21 & \multicolumn{2}{p{14.2cm}}{\todo{} Weekly Connect: show the build; collect feedback}\\
\addlinespace
Thu 22 & \todo{} 1--2 page document or 5--7 slide pitch
& \todo{} Demo video, end to end\\
\addlinespace
Fri 23 & \multicolumn{2}{p{14.2cm}}{\todo{} \textbf{First submission.}}\\
\addlinespace
Sat 24--Sun 25 & \todo{} Final combined run recorded; signed logs\newline
\todo{} File upstream issues
& \todo{} Final CI run green; video re-cut if needed\\
\addlinespace
Mon 26 & \multicolumn{2}{p{14.2cm}}{\todo{} \textbf{Final submission} --- done by 6:00 PM IST;
the deadline is 11:59 PM. Then shut down (Part 2, Step 13).}\\
\bottomrule
\end{longtable}

\subsection{Dependencies between us}

\begin{longtable}{@{}p{5.4cm}p{5.2cm}p{5.2cm}@{}}
\toprule
This & needs & Blocked meanwhile?\\
\midrule
\endhead
OpenIDConformanceTest (\Mu) & runner output (\Sh) & No --- the committed fixture is real
output shape\\
Issuer-plan runs (\Sh) & Certify environment (\Mu) & Yes, from week 4; issuer support is
built in week 3 regardless\\
Combined run (both) & Certify environment and issuer support & Yes --- the reason Certify
starts in week 3\\
GitHub Actions (\Mu) & unified compose (\Mu) and runner (\Sh) & Partly --- the unit tests
can run in CI from day one\\
Upstream issues (\Sh) & mentor answer on channel & No --- default is a GitHub issue\\
\bottomrule
\end{longtable}

\subsection{Risks, and what we do if they land}

\begin{longtable}{@{}p{5.4cm}p{11cm}@{}}
\toprule
Risk & Response\\
\midrule
\endhead
Certify cannot run locally & By Sat 10 Oct at the latest, fall back: submit Verify end to
end, the issuer support in the runner, and a documented Certify path. Per-module
conformance for one component satisfies the mandatory tasks.\\
Mukta's start slips & No word on setup progress as of 2 October. The fixture means the Java
class can start the moment setup is done; if api-test is not building by Fri 2 Oct, pair
on it that weekend.\\
Neither of us writes Java & The class is kept small and logic-free by design. Escalate by
mid-week 3 rather than absorb it quietly.\\
Screenshot mechanism fails & Only affects modules that get past DCQL. Fallback: one
documented manual upload, stated plainly rather than hidden.\\
Mentor answers change the target & Plan and variants are configuration; a version change
costs an edit and a re-run, not a rewrite.\\
\bottomrule
\end{longtable}

\end{document}
